\section{Overview \label{sec:overview}}

This section gives an overview of our approach \mainname, illustrating how it generates a well-typed program in the simply typed lambda calculus $\stlc$(a foundational functional programming language)~\cite{DBLP:journals/jsyml/Church40}.

\subsection{Problem Definition} \label{subsection:moti-problem}
We formulate the type-safe code generation problem based on two primary inputs: a language specification and a dataset of training tasks.

\newcommand{\checkBinding}[3]{\ensuremath{{#2}:{#3} \in {#1}}}
\newcommand{\abs}[3]{\ensuremath{\lambda {#1}:{#2}.\ {#3}}}
\newcommand{\vari}[1]{\ensuremath{{#1}}}
\newcommand{\snoc}[3]{\ensuremath{{#1}, {#2}:{#3}}}
\newcommand{\arr}[2]{\ensuremath{{#1}\rightarrow {#2}}}
\newcommand{\app}[2]{\ensuremath{{#1}\ {#2}}}
\newcommand{\rulename}[1]{\textsc{#1}}

\begin{figure*}[t]
\hfill
\begin{subfigure}{0.65\textwidth}
\begin{lstlisting}[basicstyle=\fontsize{8.2pt}{\baselineskip}\tt, escapeinside={@}{@}, morekeywords={data, pred, rule, func}, basewidth={0.5em, 0.2em}] 
data Prog = true | false | Var | Prog Prog | @$\tt \lambda$@Str:Type. Prog
data Type = bool | Type @$\rightarrow$@ Type      
data Context = empty | Context, Var:Type
\end{lstlisting}
\vspace{-0.5em}
\caption{The inductive data types for the syntax.} \label{figure:stlc-syntax}
\end{subfigure}
\hfill\ 

\vspace{1em}

\begin{subfigure}{0.8\textwidth}
\small

\begin{gather*}
\frac{\checkBinding{\Gamma}{x}{t}}{\stlctype{\Gamma}{\vari x}{t}} \ \ \textsc{(T-Var)} \qquad
\frac{\stlctype{\snoc{\Gamma}{x}{t_1}}{p}{t_2}}{\stlctype{\Gamma}{\left(\abs{x}{t_1}{p}\right)}{\arr{t_1}{ t_2}}} \ \ \textsc{(T-Abs)} \qquad
\frac{\stlctype{\Gamma}{p_1}{\arr{t_1}{t_2}} \quad \stlctype{\Gamma}{p_2}{t_1}}{\stlctype{\Gamma}{\app{p_1}{p_2}}{t_2}} \ \ \textsc{(T-App)} \qquad
% \frac{\stlctype{\scode{empty}}{p}{t}}{\textbf{welltyped}(p)} \ \ \textsc{(T-Entry)} 
\end{gather*}
\vspace{-0.5em}
\caption{Selected typing rules for $\stlc$. Each typing rule derives a typing judgment of the form $\stlctype{\Gamma}{p}{t}$, which asserts that program $p$ has type $t$ under context $\Gamma$} 
    \label{figure:stlc-rules}
\end{subfigure}
\vspace{-0.5em}
\caption{The definition of $\stlc$ in \mainname, comprising its syntax and some selected typing rules..
}
\label{figure:stlc-definition}
\end{figure*}

\smallskip 
\noindent \textbf{Language definition}. As shown in \autoref{figure:stlc-definition}, a language in \mainname is specified by its syntax and its type system.
%
In this paper, we define the syntax as a set of inductive data types over several built-in base types (e.g., \tcode{string}, \tcode{bool}). 
Here, the syntax of $\stlc$ consists of three data types (\autoref{figure:stlc-syntax}):
\tcode{Prog} and \tcode{Type} define programs and types in $\stlc$ as the standard simply typed lambda calculus, above the base type \tcode{bool}; and 
\tcode{Context} specifies a context for variable types, where \tcode{empty} denotes the empty context, and ${(\snoc{\Gamma}{\tt x}{\tt t})}$ extends a context $\Gamma$ by binding the type \tcode{t} to the variable \tcode{x}.

To specify the type system, we require the definition of each typing rule (which defines how types are assigned to program constructs) together with an executable type checker.
For clarity, we present our approach using the standard type notations of $\stlc$ as listed in \autoref{figure:stlc-rules}.
The type checker is expected to generate a type correctness proof (i.e., a type derivation tree, where each tree node corresponds to an application of a typing rule, like \autoref{figure:sample-type-tree}) for the typing judgement $\stlctype{\tt empty}{p}{t}$, which asserts that program $p$ has type $t$ under the empty context.

\smallskip 
\noindent \textbf{Training tasks}. Following the standard code generation paradigm, the objective is to generate correct programs from natural language specifications. Therefore, each training task in \mainname is specified by a prompt paired with the expected program.
%
\autoref{figure:sample-task} introduces an example task over $\stlc$, and \autoref{figure:sample-type-tree} illustrates the type correctness proof of the expected program generated by the type checker during type checking.

\begin{figure*}[t]
\begin{minipage}{0.3\textwidth}
    \small 
    \fbox{
        \begin{minipage}{0.98\textwidth}
\textbf{Prompt}: Implement an identity function that takes a boolean input and returns the same value; then apply it to the value true.
 
\vspace{0.5em}
\textbf{Expected program} $p^*$: 
\begin{align*}
    \app{
        (\abs{\scode{x}}{\scode{bool}}{\vari{\scode{x}}})
    }{
        \text{\scode{true}}
    }
\end{align*}
\end{minipage}
}
\caption{A training task over $\stlc$.}
\label{figure:sample-task}
\end{minipage}
\hfill
\begin{minipage}{0.65\textwidth}
\small 
\hfill\ 

\begin{prooftree}[rule margin=1ex, separation=1em]
\hypo{
  \checkBinding{
    \left\{\snoc{\scode{empty}}{\scode{x}}{\scode{bool}}\right\}}
    {\scode{x}}
    {\scode{bool}}
}
\infer1[(\textsc{T-Var})]{\stlctype{\snoc{\scode{empty}}{\scode{x}}{\scode{bool}}}{\scode{x}}{\scode{bool}}}
\infer1[(\textsc{T-Abs})]{\stlctype{\scode{empty}}{\abs{\scode{x}}{\scode{bool}}{\scode{x}}}{\arr{\scode{bool}}{\scode{bool}}}}
\hypo{~} 
\infer1[(\textsc{...})]{\stlctype{\scode{empty}}{\scode{true}}{\scode{bool}}}
\infer2[(\textsc{T-App})]{\stlctype{\scode{empty}}{\app{(\abs{\scode{x}}{\scode{bool}}{\scode{x}})}{\scode{true}}}{\scode{bool}}}
\end{prooftree}


\caption{The type correctness proof of the expected program $p^*$ in \autoref{figure:sample-task}.}
\label{figure:sample-type-tree}
\end{minipage}
\end{figure*}

\subsection{From Proof Construction to Program Synthesis} \label{subsection:moti-representation}
Given the inputs in \autoref{figure:sample-task}, \mainname aims to generate well-typed $\stlc$ programs from prompts. Our core insight is that we can construct an existential type correctness proof even when the program is unknown, thereby synthesizing the program while constructing the proof. This is possible because each typing rule restricts a part of the program structure, making the structure of the proof sufficient to determine the entire program.

For example, in \autoref{figure:sample-type-tree}, the \rulename{T-App} rule restricts the top-level operator of the program to be a function application, the \rulename{T-Abs} rule restricts the function term to be a lambda abstraction, and the application of \rulename{T-Var} defines that the body of the lambda abstraction is a variable. 
In other words, \textbf{typing rules inherently serve a program synthesis role}. This observation allows us to construct a synthesis system directly grounded in typing rules, enabling simultaneous program synthesis and type reasoning.

We illustrate the synthesis process of \mainname in detail as shown in \autoref{tab:synthesis-process}, where we use \sk{purple} to distinguish the unknown variables (i.e., placeholders to be filled in) in each goal, which will be instantiated later during the synthesis.
%
For the $\stlc$ language, \mainname begins by specifying the synthesis goal as $\stlctype{\tcode{empty}}{\sk{p}}{\sk{t}}$. This synthesis goal is a typing judgment with unknown variables, and it corresponds to constructing an existential proof of $\exists \sk{p}, \sk{t}.\; \stlctype{\tcode{empty}}{\sk{p}}{\sk{t}}$: the unknown variables in the typing judgment are exactly the existential variables to be witnessed in the proof. The initial context $\tcode{empty}$ indicates that no variables have been bound at the beginning.

During synthesis, the system incrementally applies typing rules to construct the proof, and fills in the values of $\sk{p}$ and $\sk{t}$ correspondingly. For other programming languages with different type systems, the synthesis goal can be adapted accordingly. 

\begin{table}[t]
\centering
\caption{Synthesis Process of \mainname for program in \autoref{figure:sample-task}}
\label{tab:synthesis-process}
\resizebox{\textwidth}{!}{%
\begin{tabular}{|c|l|l|l|l|l|}
\hline
\textbf{Step} & \textbf{Current Goal} & \textbf{Rule Applied} & \textbf{Substitution} & \textbf{New Subgoals} & \textbf{Current Program} \\
\hline
\multirow{2}{*}{\textbf{1}} & \multirow{2}{*}{$\stlctype{\tcode{empty}}{\sk{p}}{\sk{t}}$} & \multirow{2}{*}{\rulename{T-App}} & $\sigma_1=\{ \sk{\Gamma_1} \mapsto \tcode{empty},$ & $\stlctype{\tcode{empty}}{\sk{p_1}}{\arr{\sk{t_1}}{\sk{t_2}}}$ & \multirow{2}{*}{$\app{\sk{p_1}}{\sk{p_2}}$} \\
& & & $\sk{p} \mapsto \app{\sk{p_1}}{\sk{p_2}}, \sk{t} \mapsto \sk{t_2}\}$ & $\stlctype{\tcode{empty}}{\sk{p_2}}{\sk{t_1}}$ & \\
\hline
\multirow{2}{*}{\textbf{2}} & \multirow{2}{*}{$\stlctype{\tcode{empty}}{\sk{p_1}}{\arr{\sk{t_1}}{\sk{t_2}}}$} & \multirow{2}{*}{\rulename{T-Abs}} & $\sigma_2 = \{ \sk{\Gamma_2} \mapsto \tcode{empty},$ & \multirow{2}{*}{$\stlctype{\snoc{\tcode{empty}}{\sk{x_1}}{\sk{t_3}}}{\sk{p_3}}{\sk{t_4}}$} & \multirow{2}{*}{$\app{(\abs{\sk{x_1}}{\sk{t_3}}{\sk{p_3}})}{\sk{p_2}}$} \\
& & & $\sk{p_1} \mapsto \abs{\sk{x_1}}{\sk{t_3}}{\sk{p_3}}, \ldots\}$ & & \\
\hline
\multirow{3}{*}{\textbf{3}} & \multirow{3}{*}{$\stlctype{\snoc{\tcode{empty}}{\sk{x_1}}{\sk{t_3}}}{\sk{p_3}}{\sk{t_4}}$} & \multirow{3}{*}{\rulename{T-Var}} & $\sigma_3=\{ \sk{\Gamma_3} \mapsto (\snoc{\tcode{empty}}{\sk{x_1}}{\sk{t_3}}),$ & \multirow{3}{*}{None} & \multirow{3}{*}{$\app{(\abs{\tcode{x}}{\tcode{bool}}{\tcode{x}})}{\sk{p_2}}$} \\
& & & $\sk{p_3} \mapsto \vari{\sk{x_2}}, \sk{t_4} \mapsto \sk{t_5}\}$ & & \\
& & & $\sigma_4 = \{ \sk{x_1} \mapsto \tcode{x}, \sk{x_2} \mapsto \tcode{x}, \ldots\}$ & & \\
\hline
\multirow{2}{*}{\textbf{4}} & \multirow{2}{*}{$\stlctype{\tcode{empty}}{\sk{p_2}}{\tcode{bool}}$} & \multirow{2}{*}{\rulename{...}} & \multirow{2}{*}{$\sigma_6=\{\sk{p_2}\mapsto\tcode{true}, \ldots\}$} & \multirow{2}{*}{None } & \multirow{2}{*}{$\app{(\abs{\tcode{x}}{\tcode{bool}}{\tcode{x}})}{\tcode{true}}$} \\
& & & & & \\
\hline
\end{tabular}%
}
\end{table}

\smallskip \noindent \textbf{Step 1: applying \rulename{T-App}.}
Given the initial synthesis goal $\stlctype{\tcode{empty}}{\sk{p}}{\sk{t}}$, the LM must select a typing rule to apply. Assume the rule \rulename{T-App} is selected. The system then verifies that this application is admissible for the current goal, applies the rule, and generates the corresponding subgoals (one for each premise). More concretely, the system performs the following steps in order:
\begin{enumerate}
    \item First, \mainname $\alpha$-renames the rule by replacing all variables in the rule with fresh ones to avoid name conflicts. Under this renaming, \rulename{T-App} takes the following form.
    $$\frac{\stlctype{\sk{\Gamma_1}}{\sk{p_1}}{\arr{\sk{t_1}}{\sk{t_2}}} \quad \stlctype{\sk{\Gamma_1}}{\sk{p_2}}{\sk{t_1}}}{\stlctype{\sk{\Gamma_1}}{\app{\sk{p_1}}{\sk{p_2}}}{\sk{t_2}}} $$
    
    \item Subsequently, \mainname attempts to unify the rule's conclusion judgment with the current synthesis goal $\stlctype{\tcode{empty}}{\sk{p}}{\sk{t}}$, accumulating the induced constraints as the equations shown below. If unification fails, the rule application is considered invalid and discarded.
    $$
    \sk{\Gamma_1} = \tcode{empty} \ \ \wedge\ \ \app{\sk{p_1}}{\sk{p_2}} = \sk{p}  \ \ \wedge \ \  \sk{t_2} = \sk{t}
    $$

    \item \label{action3}Thereafter, \mainname attempts to resolve these constraints via unification, yielding a substitution that maps variables to their values. If unification fails, the rule application is rejected.
    $$\sigma_1=\big\{ \sk{\Gamma_1} \mapsto \tcode{empty}, \quad \sk{p} \mapsto \app{\sk{p_1}}{\sk{p_2}}, \quad \sk{t} \mapsto \sk{t_2}\big\}$$

    \item Finally, \mainname constructs synthesis subgoals by applying the substitution to each premise of \rulename{T-App}, which has two premises corresponding to the function and the argument of the application, 
    $\stlctype{\sk{\Gamma_1}}{\sk{p_1}}{\arr{\sk{t_1}}{\sk{t_2}}}$ and $\stlctype{\sk{\Gamma_1}}{\sk{p_2}}{\sk{t_1}}$, and \mainname will generate two subgoals:
    \begin{gather*}
    \stlctype{\tcode{empty}}{\sk{p_1}}{\arr{\sk{t_1}}{\sk{t_2}}} \quad \stlctype{\tcode{empty}}{\sk{p_2}}{\sk{t_1}}
    \end{gather*}
\end{enumerate}
Each subgoal corresponds to constructing an existential sub-proof. Specifically, the first subgoal $\stlctype{\tcode{empty}}{\sk{p_1}}{\arr{\sk{t_1}}{\sk{t_2}}}$ corresponds to proving $\exists \sk{p_1}, \sk{t_1}, \sk{t_2}.\; \stlctype{\tcode{empty}}{\sk{p_1}}{\arr{\sk{t_1}}{\sk{t_2}}}$, and the second subgoal corresponds to proving $\exists \sk{p_2}.\; \stlctype{\tcode{empty}}{\sk{p_2}}{\sk{t_1}}$. Once these sub-proofs are constructed with concrete witnesses, the substitution $\sigma_1$ determines the original existential variables: $\sk{p} = \app{\sk{p_1}}{\sk{p_2}}$ and $\sk{t} = \sk{t_2}$. \mainname proceeds to discharge the two subgoals recursively before closing the parent goal.

\smallskip \noindent \textbf{Step 2: applying \rulename{T-Abs}.}
\mainname applies the same procedure to the first subgoal $\stlctype{\tcode{empty}}{\sk{p_1}}{\arr{\sk{t_1}}{\sk{t_2}}}$. 
%
The LM will select the \rulename{T-Abs} with a high probability as this goal requires an arrow type.
\mainname then $\alpha$-renames the \rulename{T-Abs} rule:
$$
\frac{
    \stlctype{\snoc{\sk{\Gamma_2}}{\sk{x_1}}{\sk{t_3}}}{\sk{p_3}}{\sk{t_4}}
}{
    \stlctype{\sk{\Gamma_2}}{\left(\abs{\sk{x_1}}{\sk{t_3}}{\sk{p_3}}\right)}{\arr{\sk{t_3}}{\sk{t_4}}}
}
$$

By performing the remaining three actions as detailed in Step 1, \mainname again obtains a substitution and generates a new subgoal as follows.
$$\begin{array}{ll}
\textrm{Substitution: } & \sigma_2 = \big\{ \sk{\Gamma_2} \mapsto \tcode{empty}, \quad \sk{p_1} \mapsto \abs{\sk{x_1}}{\sk{t_3}}{\sk{p_3}}, \quad \sk{t_1} \mapsto \sk{t_3}, \quad \sk{t_2} \mapsto \sk{t_4}\big\}\\
\textrm{Subgoal: }& \stlctype{\snoc{\tcode{empty}}{\sk{x_1}}{\sk{t_3}}}{\sk{p_3}}{\sk{t_4}}
\end{array}$$

\smallskip \noindent \textbf{Step 3: applying \rulename{T-Var}.}
Next, \mainname handles the newly generated subgoal $\stlctype{\snoc{\tcode{empty}}{\sk{x_1}}{\sk{t_3}}}{\sk{p_3}}{\sk{t_4}}$, assuming the LM selects the \rulename{T-Var} rule. \mainname again $\alpha$-renames the rule and then obtains the substitution as shown below.
$$\begin{array}{ll}
\textrm{Rule: } & \displaystyle\frac{\checkBinding{\sk{\Gamma_3}}{\sk{x_2}}{\sk{t_5}}}{\stlctype{\sk{\Gamma_3}}{\vari{\sk{x_2}}}{\sk{t_5}}}\\[0.8em]
\textrm{Substitution: } & \sigma_3=\big\{ \sk{\Gamma_3} \mapsto (\snoc{\tcode{empty}}{\sk{x_1}}{\sk{t_3}}), \quad \sk{p_3} \mapsto \vari{\sk{x_2}}, \quad \sk{t_4} \mapsto \sk{t_5}\big\}
\end{array}$$

\rulename{T-Var} differs from previous rules in that its premise involves the membership predicate $\in$, which checks whether a variable binding exists in the context.
%
Rather than generating a new subgoal, \mainname directly invokes the predicate $\in$ to assess the validity of its arguments.
%
However, there are still unknown variables in both the premise and the synthesis goal even after applying the substitution $\sigma_3$:
$$\begin{array}{ll}
\textrm{Premise: } & \checkBinding{(\snoc{\tcode{empty}}{\sk{x_1}}{\sk{t_3}})}{\sk{x_2}}{\sk{t_5}}\\
\textrm{Goal: } & \stlctype{\snoc{\tcode{empty}}{\sk{x_1}}{\sk{t_3}}}{\sk{x_2}}{\sk{t_5}}
\end{array}$$

The unknown variables in the premise prevent \mainname from evaluating the predicate. 
%
To address this issue, \mainname will query the LM to provide an assignment to all such variables, such as follows:
%Therefore, it is necessary to instantiate these variables ($\sk{x_1}, \sk{x_2}, \sk{t_3},$ and $\sk{t_5}$ in this example) to obtain their concrete values.
%More generally, \mainname requires the LM to generate an assignment for all variables that, after applying the substitution, remain in the typing rules but are not present in any subgoals. 
%Suppose the LM generates the following assignment:
\begin{gather}
\sigma_4 = \big\{ \sk{x_1} \mapsto \tcode{x},\quad \sk{x_2} \mapsto \tcode{x},\quad \sk{t_3} \mapsto \tcode{bool},\quad \sk{t_5} \mapsto \tcode{bool} \big\} \label{formula:sig4}
\end{gather}

By further applying $\sigma_4$, the premise becomes $\checkBinding{(\snoc{\tcode{empty}}{\tcode{x}}{\tcode{bool}})}{\tcode{x}}{\tcode{bool}}$, which is completely known; then, \mainname can evaluate the predicate and confirms that it is satisfied.
%
At this time, no new subgoal is generated, so the synthesis process for this branch is complete, where the result is the composition of $\sigma_3$ and $\sigma_4$ (i.e., applying both substitutions in sequence):
\[
\sigma_4 \sigma_3 = \big\{
\sk{\Gamma_3} \mapsto \snoc{\tcode{empty}}{\tcode{x}}{\tcode{bool}},\ 
\sk{p_3} \mapsto \tcode{x},\ 
\sk{t_4} \mapsto \tcode{bool},\ 
\sk{x_1} \mapsto \tcode{x},\ 
\sk{x_2} \mapsto \tcode{x},\ 
\sk{t_3} \mapsto \tcode{bool},\ 
\sk{t_5} \mapsto \tcode{bool}
\big\}
\]

\smallskip \noindent \textbf{Step 4: returning from subgoals.}
After finishing a subgoal, \mainname will apply the obtained substitution to all remaining subgoals to synchronize the progress; and once all subgoals are closed, \mainname will compose all obtained substitutions as the final substitution for the synthesis goal.

In this example, the substitution $\sigma_4\sigma_3$ obtained in Step 3 is further composed with the substitution $\sigma_2$ from Step 2, and the resulting composition $\sigma_5 = \sigma_4\sigma_3\sigma_2$ is returned as the result of the subgoal $\stlctype{\tcode{empty}}{\sk{p_1}}{\arr{\sk{t_1}}{\sk{t_2}}}$ in Step 2. This yields a concretized and verified typing judgment, $\stlctype{\tcode{empty}}{\abs{\tcode{x}}{\tcode{bool}}{\tcode{x}}}{\arr{\tcode{bool}}{\tcode{bool}}}$.
%
Since there are two subgoals generated in Step 1, \mainname applies the composed substitution $\sigma_5$ to the second subgoal, $\stlctype{\tcode{empty}}{\sk{p_2}}{\sk{t_1}}$, yielding
\[
\stlctype{\tcode{empty}}{\sk{p_2}}{\tcode{bool}}.
\]
\mainname then applies the same procedure to this subgoal, producing the substitution 
$\sigma_6=\{\sk{p_2}\mapsto\tcode{true}\ldots\}$. 
This new substitution is then composed with the previous $\sigma_1$ and $\sigma_5$ to yield $\sigma_7=\sigma_6\sigma_5\sigma_1$, which is applied to the original
synthesis goal $\stlctype{\tcode{empty}}{\sk{p}}{\sk{t}}$, giving the final result:
\[
\stlctype{\tcode{empty}}{\app{(\abs{\texttt{x}}{\tcode{bool}}{\texttt{x}})}{\tcode{true}}}{\tcode{bool}}
\]
At this point, we have successfully constructed an existential proof of $\exists \sk{p}, \sk{t}.\; \stlctype{\tcode{empty}}{\sk{p}}{\sk{t}}$, with the existential variables instantiated as $\sk{p} = \app{(\abs{\texttt{x}}{\tcode{bool}}{\texttt{x}})}{\tcode{true}}$ and $\sk{t} = \tcode{bool}$. The synthesized program $\sk{p}$ is thus guaranteed to be well-typed.

Performing program synthesis via existential proof construction embodies our two key principles: \textbf{Type Explicitness} and \textbf{Context Locality}.
\begin{itemize}
    \item \textbf{Type Explicitness.} Each synthesis step corresponds to applying a typing rule, which directly mirrors the structure of the type correctness proof. As shown in the example, the sequence of rule applications (\rulename{T-App}, \rulename{T-Abs}, \rulename{T-Var}, etc.) constructs a complete proof tree, making the type reasoning explicit throughout the synthesis process.
    \item \textbf{Context Locality.} At each step, the synthesis goal is a partially instantiated typing judgment that contains the current typing context. For instance, in Step 2, the goal $\stlctype{\tcode{empty}}{\sk{p_1}}{\arr{\sk{t_1}}{\sk{t_2}}}$ explicitly provides the context $\tcode{empty}$ and indicates that the program $\sk{p_1}$ to be synthesized must have an arrow type $\arr{\sk{t_1}}{\sk{t_2}}$. This local type information guides the LM to select the \rulename{T-Abs} rule, which produces a lambda abstraction.
\end{itemize}